%=====================================================================
\section{Embedding-Space Analysis}
\label{sec:analysis}
%=====================================================================

Why does \cinder{} find concepts that \ember{} misses?
We analyze the features before any edit (\expnum{38} concepts, 15 added after the method was frozen; 5 seeds), judge-free wherever possible (\S\ref{app:analyses}).

\paragraph{Sharpening makes the concept pure and valuable}
We seed a sparse feature on each concept's lexicon tokens and run it to its fixed point (Figure~\ref{fig:mechanism}b).
In the raw matrix it drifts toward other rare tokens and ends up only \expnum{24\%} concept-only, although it is among the most valuable features; whitening makes it pure (\expnum{86\%}) but ordinary (\expnum{70th} percentile of feature energies), and sharpening keeps it pure and makes it valuable (\expnum{91st}, higher in \expnum{37 of 38} concepts).
In full fits, a lexicon-matching concept feature appears in \expnum{56\%} of raw, \expnum{49\%} of whitened-only and \expnum{69\%} of sharpened runs, with fewer neutral tokens (\expnum{3.0} vs.\ \expnum{4.5}).
Its nearest neighbour becomes a related topic (magic for Harry Potter) rather than a feature of unrelated rare tokens, and \expnum{97\%} of the raw feature's excess squared norm lies along $\hat\mu_B$ (\S\ref{app:geometry}).

\paragraph{Widths follow the data}
On the sharpened matrix, \ember{}'s 1\% cut stops before the natural tail for \expnum{25 of 38} concepts, and \gs{} widths follow the tail (Spearman \expnum{0.67} vs.\ \expnum{0.27}; mean error \expnum{30} vs.\ \expnum{42} tokens, $p < 0.001$), though they compress its range; \gs{} changes which tokens a feature keeps, not its direction (cosine \expnum{0.97}), and helps only after \cs{} (\S\ref{app:gs}).

\paragraph{The two-way judge recovers concepts lost in the summary}
On 19 concepts, \twj{} accepts the concept or a broader version of it in \expnum{82 of 91} runs, against \expnum{61} for \ember{}'s judge, with \expnum{1 of 110} acceptances unrelated, though most added features are broader than the concept.
On 15 concepts added later, the lexicon path generalizes poorly, whereas showing the judge each token's source word does (TJ with source words \expnum{49}, TJ \expnum{33} of 54 runs; \S\ref{app:twj}).

\paragraph{The gain is not circular}
The concept's page, our sharpening target, contains the sentences whose tokens we score, yet sharpening toward page sentences outside the sample or toward other pages about the concept raises the strict purity of the concept feature as much (\expnum{$+0.06$} vs.\ \expnum{$+0.05$}; Holm $p = \expnum{0.034}$, \expnum{0.081}; pre-registered), while another concept's page destroys it (found in \expnum{2 of 54} runs; \S\ref{app:cs-circ}).

\paragraph{Abstract concepts stay out of reach}
On 13 abstract concepts such as Stoicism and Gothic architecture, the top-ranked feature is specific to the concept in \expnum{0 of 65} runs under both \ember{} and \cs{}, against \expnum{13} and \expnum{23 of 40} for entity concepts; their top features are parent topics and discourse words, perhaps a register they share (untested).
Dilution with adjacent pages hurts \ember{} more than \cs{} (\S\ref{app:robustness}), and features extrapolate only weakly beyond the sample (\S\ref{app:extrapolation}).

\paragraph{The first layer reads the shared mean}
The model's first layer reads \ember{}'s and \cinderz{}'s edits at about twice a random direction, \cinder{}'s at \expnum{1.2$\times$} (Figure~\ref{fig:read}, \S\ref{app:edit}).
The gap is the mean term, read at \expnum{2.9$\times$} a random direction on Gemma and \expnum{20$\times$} on Llama: with $\mu_B$ projected out, \ember{}'s and \cinder{}'s edits are read alike (\expnum{1.31$\times$} vs.\ \expnum{1.32$\times$}).
